
\documentclass[11pt,letterpaper]{article}

\usepackage[margin=1in]{geometry}
\usepackage{titlesec}
\usepackage[hidelinks]{hyperref}
\usepackage{fancyhdr}
\usepackage{booktabs}

% Page formatting
\pagestyle{fancy}
\fancyhf{}
\fancyfoot[R]{\footnotesize\bfseries Page \thepage}
\renewcommand{\headrulewidth}{0pt}
\renewcommand{\footrulewidth}{0pt}

% Paragraph formatting
\setlength{\parindent}{1em}
\setlength{\parskip}{0pt}

% Section formatting
\titleformat{\section}
  {\normalfont\normalsize\bfseries}
  {}{0pt}{\MakeUppercase}
\titlespacing*{\section}
  {0pt}{16pt}{6pt}

\begin{document}

\noindent\hfill
{\footnotesize\bfseries CSCI 441 Neural Networks, Project Update 2}
\par

\vspace{1.5em}

\begin{center}
    {\Large\bfseries
    Multi-Class Traffic Sign Classification Using
    \\Convolutional Neural Networks\par}

    \vspace{0.5em}

\end{center}

\vspace{1em}

\noindent{\Large\scshape Experiment Design\par}

\vspace{1em}

\section*{Research Questions}

\begin{itemize}
    \item \textbf{RQ1:} How do different data augmentation strategies affect CNN validation accuracy on clean and altered traffic-sign images?
    \item \textbf{RQ2:} How does increasing the number of convolutional layers in a CNN affect validation accuracy on clean and altered traffic-sign images?
\end{itemize}

\section*{Dataset}
This project uses the German Traffic Sign Recognition Benchmark (GTSRB), containing 51,839 images across 43 traffic-sign classes, with 39,209 training and 12,630 test images \cite{stallkamp2011gtsrb}.
Images vary in resolution, illumination, and viewing angle.
The official test set is held out until final evaluation.

\smallskip

\noindent\textbf{Dataset: }\par
\noindent\url{https://www.kaggle.com/datasets/meowmeowmeowmeowmeow/gtsrb-german-traffic-sign}

\section*{Data Preprocessing}
Images are converted to RGB, resized to $48 \times 48$ pixels, and normalized to $[0,1]$. 
Training-only augmentation includes random rotation ($\pm 10^\circ$), brightness adjustment ($\pm 0.1$ on normalized pixels), and zoom ($0.9\times$ to $1.1\times$). 

\section*{Experimental Variables}

\noindent\textbf{RQ1:}
The training augmentation strategy is the independent variable.
All configurations use the baseline architecture, with fixed training hyperparameters and dataset splits.\\

\noindent\textbf{RQ2:}
Convolutional depth is the independent variable.
Channel width, kernel size, training hyperparameters, and dataset splits remain fixed, with no training augmentation.
Parameter count and pooling placement also vary with depth, a limitation.\\

\noindent For both research questions, evaluation condition is a second independent variable and validation accuracy is the dependent variable.

\clearpage

\section*{Models and Algorithms}
Custom CNNs with two, three, and four convolutional layers are implemented using TensorFlow/Keras.
Each convolutional block uses 64 filters, batch normalization, and ReLU activation.
All architectures use two max-pooling operations, global average pooling, and a shared classification head.

\section*{Baseline}

A two-layer CNN trained without augmentation is the baseline for both research questions because it is the reference level of both independent variables: the shallowest depth and no augmentation.

\section*{Training Configuration}

\begin{center}
\begin{tabular*}{\textwidth}{@{\extracolsep{\fill}}ll@{}}
\toprule
\textbf{Hyperparameter} & \textbf{Value} \\
\midrule
Filters per convolutional layer & 64 \\
Convolutional kernel size & $3 \times 3$ \\
Max-pooling size & $2 \times 2$ \\
Dropout & 0.2 \\
Dense units & 128 \\
Epochs & 42 \\
Batch size & 64 \\
Learning rate & 0.001 \\
Optimizer & Adam \\
Loss & Sparse categorical cross-entropy \\
Random seed & 42 \\
\bottomrule
\end{tabular*}
\end{center}

\section*{Experiment Procedure}

\begin{enumerate}
    \item Split the training set with a fixed, class-stratified, track-aware split, holding out about 20\% of each class's tracks for validation.\footnote{
        \href{https://github.com/joolaoye/csci441-project/blob/main/configs/splits/gtsrb_v1.json}
        {GTSRB split manifest.}
    }

    \item Train five augmentation configurations (none, rotation, zoom, brightness, and combined) for RQ1 and three depth configurations for RQ2, sharing the baseline.

    \item Select each model's checkpoint with the highest clean validation accuracy.

    \item Evaluate each selected checkpoint under five conditions: clean, $10^\circ$ rotation, $1.1\times$ zoom, brightness adjustment of $+0.1$, and downsampling to $24 \times 24$  then upsampling to $48 \times 48$.
    No model trains on downsampling.

    \item Evaluate the selected checkpoints once on the clean test set.
\end{enumerate}

\section*{Evaluation Metrics}
Validation accuracy, the proportion of correctly classified validation images, is the primary metric. Macro-F1 is secondary because classes are imbalanced.

\clearpage

\section*{Result Analysis}
Accuracy and macro-F1 will be compared across configurations and evaluation conditions, relative to the baseline and the clean condition, using tables and plots.
With one run per configuration, small performance differences will be interpreted cautiously.
Learning curves and confusion matrices will support interpretation of convergence and classification errors.

\bibliographystyle{IEEEtran}
\bibliography{references}

\end{document}
